% TOP_PROBLEM: {"problemId": "problem.kaplansky-s-direct-finiteness-conjecture", "canonicalTitle": "Kaplansky's Direct Finiteness Conjecture", "releaseRank": 440, "primaryDomain": "algebra_representation_category", "primaryDomainLabel": "Algebra, representation & category theory", "displayStatus": "open_with_solved_subcases", "releaseStatus": "open_with_solved_subcases"}
% !TeX program = lualatex
% Definition and sources only; this file is not an LLM solution attempt.
\documentclass[11pt]{article}
\usepackage[margin=25mm]{geometry}
\usepackage{fontspec}
\setmainfont{Latin Modern Roman}
\usepackage{amsmath,amssymb,mathtools}
\usepackage{unicode-math}
\setmathfont{Latin Modern Math}
\usepackage{xurl,hyperref}
\hypersetup{colorlinks=true,urlcolor=blue}
\setcounter{secnumdepth}{0}
\setlength{\parindent}{0pt}
\setlength{\parskip}{5pt}
\setlength{\emergencystretch}{3em}
\begin{document}
\section*{\texorpdfstring{Kaplansky's Direct Finiteness Conjecture}{Kaplansky's Direct Finiteness Conjecture}}\label{440}
\subsection{Definitions and mathematical statement}
For a group $G$ and field $K$, the group ring $K[G]$ consists of finite sums $\sum_g a_g g$, with coefficientwise addition and multiplication extended bilinearly from the group law. Its identity is $1_K e_G$. A unital ring is directly finite if a right inverse is always a left inverse. The assertion is
\[
 \forall G\ \forall K\ \forall a,b\in K[G],
 \qquad ab=1\Longrightarrow ba=1.
\]
The group may have torsion and the field may have any characteristic. Stable finiteness means direct finiteness of every matrix ring $M_n(K[G])$. The catalogue records a known equivalence of the universal direct- and stable-finiteness targets; the displayed problem itself is the direct-finiteness assertion, with no bound on the supports of $a,b$.
\par\smallskip\textit{Formulation and concept references:} \href{https://web.ma.utexas.edu/users/juschenko/files/soficgroups.pdf}{[S1]}.

\subsection{Short English statement}
In every group ring over a field, must reversing a product that equals one still give one?
\subsection{Sources}
\begin{itemize}
\item[S1]Kate Juschenko, Sofic Groups, Chapter 5: Some conjectures that are valid for sofic groups.
\url{https://web.ma.utexas.edu/users/juschenko/files/soficgroups.pdf}.
\textit{Formulation source.}
\end{itemize}
\end{document}
